\cotefile{46}
\scanpage{81}
% archive sheet number: 80
% block text | lines 1-2
Soit $\Gamma$ un groupe, $G$ un $\Gamma$-groupe, $M$ un $\Gamma$-\deleted{\unsure{groupe abélien}}\inserted{ensemble}. Supposons que $G$ opère sur $M$, et qu'on ait
% block display | line 3
\[
\unsure{\gamma|\sigma}(gm)=(\gamma g)(\gamma m)
\]
% block text | line 4
Soit alors $f$ un $1$-cocycle de $\Gamma$ à valeurs dans $G$
% block display | line 5
\[
f(\gamma\gamma')=f(\gamma)\cdot\gamma f(\gamma')
\]
% block text | lines 6-8
La donnée de $f$ permet de ``tordre'' $M$, de façon \unsure{précise} de définir une autre façon d'opérer de $\Gamma$ sur $M$, savoir
% block display | line 8
\[
\overset{\deleted{\unsure{\sim}}}{\gamma}_f m=f(\gamma)\cdot\gamma\cdot m
\]
% block display | lines 9-11
\[
\left[\begin{aligned}
\text{Car }\quad\gamma_f(\overset{\deleted{\unsure{\sim}}}{\gamma}'_f m)&=f(\gamma)\gamma(f(\gamma')\gamma'm)=f(\gamma)\gamma f(\gamma')\,\gamma\gamma'm\\
&=f(\gamma\gamma')\gamma\gamma'm=\overset{\deleted{\unsure{\sim}}}{(\gamma\gamma')}_f m\\
e_f m&=f(e)e\cdot m=m
\end{aligned}\right]
\]
% block text | lines 12-17
On désignera par $M_f$ le nouveau \deleted{\unsure{module}}\inserted{$\Gamma$-ensemble} ainsi construit; en fait, on ne s'est pas servi de la loi de groupe de $M$; en particulier, faisant opérer $G$ dans lui-même par automorphismes intérieurs, on \deleted{\unsure{trouve sur $G_f$ un module|a sur $G_f$ un module}} voit que
% block display | line 18
\[
\gamma(\sigma(g)\cdot g')=\sigma(\gamma g)\,\gamma g'
\]
% block text | lines 19-20
car le \deleted{\unsure{premier}}\inserted{premier} \unsure{mbr|mbre} est $\gamma(gg'g^{-1})=\gamma g\,\gamma g'(\gamma g)^{-1}$, c'est à dire $=$ le second membre. Ainsi, posant
% block display | line 21
\[
\gamma_f g=f(\gamma)\gamma g f(\gamma)^{-1}
\]
% block text | line 22
on obtient sur $G$ une nouvelle loi de $\Gamma$-groupe.

% block text | line 23
Notons \unsure{enfin} \unsure{.|:} nous avons
% block display | line 24
\[
\gamma_f(gm)=\gamma_f(g)\,\gamma_f(m)
\]
% block text | line 25
car le \deleted{deuxième}\inserted{premier} mbr est
% block display | lines 26-27
\[
\begin{aligned}
&f(\gamma)\gamma(gm)=f(\gamma)\gamma(g)\,\gamma(m)\\
&f(\gamma)\gamma g f(\gamma)^{-1}f(\gamma)\gamma(m)
\end{aligned}
\]
% block text | lines 26-27 right
et le deuxième \unsure{i.e.} pareil.
